\scanpage{7}
% archive foliation: 6
\uline{Lemme} Supposons les $L_n$ et $\pi_n(L_*)$
projectifs de \unsure{t.f.|r.f.} sur $k$, et supposons $k$
commutatif à \uline{spec connexe} (donc
$\forall M$ projectif sur $k$, $\operatorname{rg}_k M\in\mathbf N$ est défini).
\margin{\unsure{on peut aussi} \illegible{1} \unsure{au sens} \illegible{1} \unsure{projectifs}}

Conditions équivalentes

a) $L^!_*$ à degrés bornés

b) \deleted{\unsure{$L_*$}} Des $M_i$, $N_j$ \unsure{n.} nuls sauf
un nb fini t.q. $L_*\simeq$ somme \uline{finie}
de modules $k$-simpliciaux $\Lambda^i\Phi_*\otimes M_i$ $(i\geq1)$,
$\Lambda^j\Psi_*\otimes N_j$ $(j\geq1)$ ($M_i$ et $N_j$ proj. de \unsure{t.f.|r.f.})

b') (Si $k$ principal, plus génér. si tout
module proj de \unsure{t.f.|r.f.} sur $k$ est libre).
$L_*\simeq$ somme \uline{finie} de modules $k$-
simpliciaux $\Lambda^i\Phi_*$ \inserted{$(i\geq1)$}, $\Lambda^j\Psi_*$ $(j\geq0)$

c) $\operatorname{rang}_k L_n$ est une fonction
polynomiale en $n$ pour $n$ grand

d) $\exists N$ \inserted{et $c>0$} tel que pour $n$ assez
grand, on ait
\[
\operatorname{rang}_k L_n\leq c n^N
\]

[\deleted{\unsure{Cor.}} De plus, \deleted{\unsure{p}} pour $N$ fixé, la condition
que l'on ait $\operatorname{rang}_k L_n=O(n^N)$
signifie que $M_i$, $N_j$ sont nuls pour $i,j>N$;
\deleted{\unsure{et si $N$ est le plus petit possible} (\unsure{donc} $M_N+N_N\neq0$)}
\unsure{cf. b) ci-dessus} \deleted{\illegible{6}}
(\unsure{voir 1 et 2})]
\margin{\unsure{et si} \illegible{2}
$\displaystyle\unsure{\lim_{n\to\infty}\frac{\operatorname{rg}_k L_n}{n^N}
=\mu_N+\nu_N}$}

\uline{Dém.} \deleted{Équivalence des \unsure{conditions}}
\unsure{ci-dessus}), $\mathrm b)\Longleftrightarrow\mathrm b')$ trivial,
\deleted{\unsure{a)\ensuremath{\Rightarrow}}}
$\mathrm c)\Rightarrow\mathrm d)$ triv., $\mathrm a)\Rightarrow\mathrm c)$ et $\mathrm d)\Rightarrow\mathrm a)$
immédiat par les formules explicites
